% =====================================================================
% CHAPTER 1 — INTRODUCTION (~8% of total length)
%
% Job: by the end, the reader believes (a) the insufficient-statistic
% claim is a real, testable question, and (b) trading P/L is the right
% test of it. Everything else hangs off that.
%
% Structure:
%   1.1 Motivation — pipelines compress text to a scalar; compression
%       destroys information (data-processing inequality); is what's
%       destroyed PRICED? That one-paragraph arc is the whole thesis.
%   1.2 Research questions — RQ1–RQ4 below mirror layers L1–L4 of the
%       research plan. Keep the numbering stable: Chapters 4–7 each
%       answer one, and the conclusion closes them one by one.
%   1.3 Contributions
%   1.4 Scope & limitations
%   1.5 Outline
%
% PITFALL: don't oversell L1. It is the supervisor's baseline and the
% testbed — the novelty lives in L2 (consensus→crowding) and L3
% (reliability→sizing). Frame it that way from page 1.
% =====================================================================
\chapter{Introduction}
\label{ch:introduction}

\section{Motivation}
\label{sec:motivation}

\todo{Open with the compression arc: news $\to$ one number $\to$ trade. What did the one number throw away?}
A sentiment trading pipeline compresses each day's news into a single
scalar $S_t$ and trades on it. When multiple models, prompts, and
stochastic samples score the same text, they produce a
\emph{distribution} of readings, and distinct properties of that
distribution --- the mean, the cross-model agreement, the measurement
reliability --- may carry distinct, economically meaningful
information that the scalar destroys. [Develop in 2--3 paragraphs:
(i) the rise of \gls{llm} sentiment pipelines and the industry's
convergence on a few foundation models; (ii) why machine consensus may
proxy for \emph{crowded positioning} rather than confidence; (iii) why
profit-and-loss, not classification accuracy, is the decisive test.]

\section{Research Questions}
\label{sec:research-questions}

% These are the pre-registered RQs from the research plan — keep the
% wording aligned with what was registered; deviations invite the
% "did you move the goalposts?" question.
This dissertation addresses four questions, one per analysis layer:
\begin{enumerate}
  \item \textbf{RQ1 (L1, baseline):} Can a daily aggregated sentiment
        signal, traded via a tuned threshold rule, generate profit on
        index futures against price-only and buy-and-hold baselines?
  \item \textbf{RQ2 (L2, consensus $\to$ crowding):} Does the level of
        agreement among independent models about a news item predict
        how its initial price reaction resolves --- reversal after
        high-consensus readings versus drift after low-consensus
        readings?
  \item \textbf{RQ3 (L3, reliability $\to$ sizing):} Does a trading
        rule that knows the measurement error of its sentiment signal
        --- shrinking unreliable scores and sizing positions by
        signal-to-noise --- outperform the fixed-threshold rule on the
        same signal?
  \item \textbf{RQ4 (L4, conditional):} Do sentiment scorers shift
        systematically on machine-written versions of the same news?
        % If the L4 gate closed, demote this to a future-work
        % paragraph in the conclusion and delete Chapter 7.
\end{enumerate}

\section{Contributions}
\label{sec:contributions}

% Each bullet should map to evidence in a specific chapter.
The main contributions are:
\begin{itemize}
  \item [A leakage-controlled replication of the threshold strategy of
        \textcite{mitic2026llm} on index futures (\cref{ch:layer1}).]
  \item [The first event-study evidence on whether cross-model
        sentiment consensus predicts reversal vs.\ drift
        (\cref{ch:layer2}).]
  \item [The first generalizability-theory variance decomposition of
        \gls{llm} sentiment measurement in finance, and a test of its
        economic value via reliability-weighted position sizing
        (\cref{ch:layer3}).]
  \item [\dots]
\end{itemize}

\section{Scope and Limitations}
\label{sec:scope}

[Declare upfront: daily horizon; FTSE/DOW/Hang Seng futures plus
index-constituent single names for the event study; news sources and
period; no intraday execution; transaction costs as sensitivity, not
microstructure realism. Also state the publishable-null stance: a
clean negative on RQ2 or RQ3 is a finding about machine-read
efficiency, not a failure.]

\section{Outline}
\label{sec:outline}

\Cref{ch:literature} reviews the four literatures this work draws on.
\Cref{ch:data} describes data, pipeline, and the rigor protocol.
\Cref{ch:layer1,ch:layer2,ch:layer3} present the baseline strategy,
the consensus--crowding study, and the reliability--sizing study
respectively. [\Cref{ch:layer4} reports the echo experiment ---
delete if gated.] \Cref{ch:discussion} interprets the results and
\cref{ch:conclusion} concludes.
